% group2_quantum_chemistry_synthesis.tex
% GeoVac Group 2 Synthesis: Quantum Chemistry / Natural Geometry Hierarchy Arc
% Status: Project archive synthesis, June 2026
% Audience: quantum chemists, electronic-structure theorists

\documentclass[aps,pra,twocolumn,superscriptaddress,longbibliography]{revtex4-2}

\usepackage{amsmath,amssymb,amsthm,graphicx,xcolor,bm,braket,booktabs}
\usepackage{hyperref}

\newtheorem{theorem}{Theorem}
\newtheorem{proposition}{Proposition}
\newtheorem{corollary}{Corollary}
\newtheorem{observation}{Observation}
\newtheorem{conjecture}{Conjecture}

\begin{document}

\title{The Natural Geometry Hierarchy:\\
A Synthesis of the GeoVac Quantum-Chemistry Arc}
\thanks{Group~2 synthesis in the Geometric Vacuum series}

\author{J.~Loutey}
\affiliation{Independent Researcher, Kent, Washington}

\date{June 26, 2026}

%% ========================================================================
%% ABSTRACT
%% ========================================================================

\begin{abstract}
The Geometric Vacuum (GeoVac) framework discretizes non-relativistic
electronic structure on \emph{natural geometries}:\ for each separable
quantum system there is a coordinate system in which the Coulomb
singularities become boundary conditions and the Laplace--Beltrami
operator discretizes on a sparse graph with quantum-number labels for
nodes.  This synthesis assembles the framework's quantum-chemistry
arc---the nine Group~2 papers spanning one-electron atoms, the
two-center one-electron ion, the two-electron atom, the two-center
two-electron molecule, and the core--valence and polyatomic
extensions---into a single coherent narrative organized by the
\emph{natural geometry hierarchy}.  The hierarchy is a grid indexed by
the number of nuclear centers and the number of electrons:\ $S^3$
(Fock) for atoms~\cite{loutey_paper7}, the prolate spheroid for
$\mathrm{H}_2^+$~\cite{loutey_paper11}, hyperspherical coordinates for
He~\cite{loutey_paper13}, molecule-frame hyperspherical coordinates for
$\mathrm{H}_2$~\cite{loutey_paper15}, and composed fiber bundles for
core--valence and polyatomic systems~\cite{loutey_paper17,
loutey_paper19}.  Three threads run through the arc.  \emph{First}, an
algebraic-first program:\ the angular content is computed exactly from
Gaunt coefficients and Wigner $3j$ symbols at every level, the
prolate-spheroidal electron--electron repulsion is replaced by an
algebraic Neumann expansion~\cite{loutey_paper12}, the radial solvers
reduce to three-term Laguerre recurrences, and the split-region
Legendre nuclear coupling terminates exactly via the $3j$ triangle
inequality~\cite{loutey_paper15}; what currently survives as quadrature
($Z_{\mathrm{eff}}$ screening, the adiabatic potential curves, the
hyperradial coupling) is named explicitly, and finding algebraic
replacements is an ongoing goal rather than a deficiency.
\emph{Second}, two proven structural negative theorems delimit the
design space:\ in a single-center molecular Sturmian basis with a
shared momentum scale the one-electron Hamiltonian and overlap obey the
two-term identity $H_{ij} = -\tfrac12 p_0^2 S_{ij} + (1-\beta_j)V_{ij}$
and share a single orthogonal $SO(4)$ Wigner-$D$ congruence, so the
generalized eigenvalues are independent of the internuclear distance and
cannot bind unless the Sturmian scale $\beta_k(R)$ itself becomes
$R$-dependent~\cite{loutey_paper8}; and a graph
built by concatenating atom-centered lattices has an $R$-independent
intra-atomic kinetic energy, producing a monotonically attractive
potential-energy surface with no equilibrium~\cite{loutey_fci_mol}.
\emph{Third}, honest accuracy ceilings:\ the framework reaches
$0.0002\%$ for $\mathrm{H}_2^+$, $0.022\%$ (properly variational, raw
$l_{\max}=7$) down to $0.004\%$ (a non-variational cusp extrapolation)
for He, $96.0\%$ of $D_e$ for $\mathrm{H}_2$, and---in the
composed architecture, whose accuracy bottleneck is the
Phillips--Kleinman pseudopotential---$5.3\%$, $11.7\%$, and $19.4\%$
equilibrium-bond-length error for LiH, $\mathrm{BeH}_2$, and
$\mathrm{H}_2\mathrm{O}$ respectively.  The classical-solver
investigation is complete:\ its accuracy ceilings are characterized,
and the framework is positioned as a zero-parameter, structurally
sparse \emph{research instrument}---a computationally principled
alternative whose primary downstream value (qubit-Hamiltonian sparsity
for quantum simulation) lives in the quantum-computing arc---not as a
replacement for production quantum chemistry.  The discrete graph and
the continuous Laplace--Beltrami operator are dual descriptions
connected by proven conformal maps; the interpretive question of which
is more fundamental is left open.
\end{abstract}

\maketitle

%% ========================================================================
\section{Introduction}
\label{sec:intro}
%% ========================================================================

The Group~3 foundations arc of GeoVac establishes that the hydrogen
atom's momentum-space Schr\"odinger equation maps, by Fock's 1935
stereographic projection~\cite{fock1935}, onto the free Laplacian on
the unit three-sphere $S^3$, and that a discrete graph on
quantum-number labels converges to the $S^3$ Laplace--Beltrami operator
in the continuum limit~\cite{loutey_paper0, loutey_paper1,
loutey_paper7}.  That arc answers a question about a single
one-electron atom.  The present synthesis takes up the question the
foundations arc opens but does not close:\ what is the right
discretization for a \emph{molecule}, or for an atom with \emph{more
than one electron}?

The organizing answer is the \emph{natural geometry principle}, stated
first in Paper~11~\cite{loutey_paper11}:\ every separable quantum
system has a natural lattice, namely the lattice that discretizes the
Laplace--Beltrami operator in the coordinate system where the problem
separates.  For a one-center one-electron system that coordinate system
is $S^3$; for a two-center one-electron system it is the prolate
spheroid; for a one-center two-electron system it is the hyperspherical
coordinate system; for a two-center two-electron system it is the
molecule-frame hyperspherical system; and for systems with electrons on
separated length scales (a tight core and a diffuse valence shell) the
natural object is a \emph{composed} geometry---a fiber bundle in which
each electron group inhabits its own natural coordinate system, coupled
by screening functions and pseudopotentials.  Crucially, this is not a
linear tower in which each level subsumes the previous one:\ the
Level-4 two-electron coordinates $(R_e, \alpha)$ are simply undefined
for a single electron, so the hierarchy is a two-dimensional grid
indexed by (number of centers) $\times$ (number of
electrons)~\cite{loutey_paper17}.

This grid structure is the reason the chemistry arc is not a single
algorithm but a sequence of natural geometries, each with its own
separation of variables, its own algebraic content, and its own
accuracy ceiling.  Table~\ref{tab:hierarchy} summarizes the hierarchy
and the headline result at each level; the body of this synthesis
unpacks each row, the algebraic structure that makes it work, and the
two proven theorems that mark the boundaries of what natural geometry
can and cannot do.

\begin{table*}[t]
\centering
\caption{The GeoVac natural geometry hierarchy for quantum chemistry.
Each level is the coordinate system in which the corresponding system
separates; the result column reports the headline figure at its stated
tier, with the source paper.  Composed (Level~5) results are limited by
the Phillips--Kleinman pseudopotential, the characterized accuracy
ceiling of the composed architecture (Sec.~\ref{sec:composed}).}
\label{tab:hierarchy}
\begin{tabular}{l l l l c}
\toprule
Level & System (centers, electrons) & Natural geometry & Headline result & Paper \\
\midrule
1 & atom (1, 1) & $S^3$ (Fock projection) & $<0.1\%$ (one-electron) & \cite{loutey_paper7} \\
2 & $\mathrm{H}_2^+$ (2, 1) & prolate spheroid & $0.0002\%$ energy (spectral Laguerre) & \cite{loutey_paper11} \\
3 & He (1, 2) & hyperspherical & $0.022\%$ (2D var.\ raw) to $0.004\%$ (cusp-extrapolated, non-variational); $0.19\%$ (graph-native CI) & \cite{loutey_paper13} \\
4 & $\mathrm{H}_2$ (2, 2) & molecule-frame hyperspherical & $96.0\%$ of $D_e$ ($l_{\max}=6$, 61 channels) & \cite{loutey_paper15} \\
4N & LiH (2, 4) & full mol-frame ($SO(12)$) & $R_{\mathrm{eq}}\approx 63.5\%$ ($l_{\max}=2$, 2D variational; unbound $D_e$, no PK) & \cite{loutey_paper17} \\
5 & LiH (core+valence) & composed (L3 $\oplus$ L4) & $R_{\mathrm{eq}}$ $5.3\%$ ($l$-dependent PK) & \cite{loutey_paper17} \\
5 & $\mathrm{BeH}_2$ (polyatomic) & composed $+$ 1-RDM exchange & $R_{\mathrm{eq}}$ $11.7\%$ & \cite{loutey_paper17} \\
5 & $\mathrm{H}_2\mathrm{O}$ (triatomic) & composed (five-block, uncoupled) & $R_{\mathrm{eq}}$ $19.4\%$ & \cite{loutey_paper17} \\
\bottomrule
\end{tabular}
\end{table*}

GeoVac is presented throughout as a discretization framework that
exploits natural geometry, not as a new foundation for quantum
mechanics.  The discrete graph and the continuous Laplace--Beltrami
operator are dual descriptions connected by the proven conformal maps
of the foundations arc, and the interpretive question of ontological
priority is left to the reader.  The framework's concrete
advantages---a zero-parameter construction from nuclear charges and
geometry alone, angular-momentum selection rules baked into the basis,
exact algebraic matrix elements, and structural
sparsity---are independent of that interpretive question and are the
properties this synthesis foregrounds.

%% ========================================================================
\section{Level 1--2: From the Atom to the Diatomic Ion}
\label{sec:level12}
%% ========================================================================

\subsection{The atomic baseline and the bond sphere}

The one-electron atomic baseline is the $S^3$ graph of the foundations
arc:\ the discrete Fock lattice reproduces the Rydberg spectrum and
solves one-electron atoms to better than $0.1\%$ with no free
parameters~\cite{loutey_paper1, loutey_paper7}.  The natural first
question for chemistry is whether the same $S^3$ object handles a
\emph{molecule}.  Paper~8~\cite{loutey_paper8} pursues this directly:\
it shows that a diatomic molecule is geometrically a single $S^3$ with
two weighted poles---the \emph{bond sphere}---in which the internuclear
distance $R$ encodes as a polar angle $\gamma$ through
$\cos\gamma = (p_0^2 - p_R^2)/(p_0^2 + p_R^2)$, and all cross-atom
matrix elements are $SO(4)$ Wigner $D$-matrix elements evaluated at the
bond rotation $e^{i\gamma A_y}$.  The geometry is consistent, and it
yields exact selection rules for $\sigma$-bonds:\ the bond rotation
preserves $m=0$ within $l=1$ and does not mix $s$ and $p$ states,
$D^2_{10,10}(\gamma) = 1$ and $D^2_{00,10}(\gamma) = 0$ for all
$\gamma$~\cite{loutey_paper8}.

\subsection{The Sturmian structural theorem (a guardrail)}

The bond sphere construction also produces the arc's first proven
\emph{negative} theorem, and it is load-bearing for everything that
follows.  In the molecular Sturmian basis~\cite{avery2004} with a shared momentum scale
$p_0$, the one-electron Hamiltonian matrix is
\begin{equation}
H_{ij} = -\frac{p_0^2}{2}\,S_{ij} + (1-\beta_j)\,V_{ij},
\label{eq:sturmian_structural}
\end{equation}
where $S_{ij}$ is the overlap matrix, $V_{ij}$ the nuclear-attraction
matrix, and $\beta_j$ is the Sturmian eigenvalue of orbital~$j$.  The
Hamiltonian is \emph{not} a scalar multiple of $S$:\ both $(H,S)$ carry
their entire $R$-dependence through the same $SO(4)$ Wigner-$D$
congruence $D^{(n)}(\gamma(R))$, so the generalized eigenvalues
$\varepsilon_k = \tfrac{1}{2}p_0^2 - p_0^2/\beta_k$ (which equals
$-p_0^2/2$ only in the degenerate case $\beta_k=1$) are independent of
the bond angle $\gamma$, hence independent of the internuclear
distance~$R$~\cite{loutey_paper8}.  A basis that produces
$R$-independent eigenvalues cannot describe binding, which is by
definition an $R$-dependent lowering of the energy.  Paper~8 sharpens
this further:\ for a heteronuclear diatomic with charge ratio
$Z_A/Z_B > 1$ there is no atom-centered shared-$p_0$ Sturmian basis
that binds both centers simultaneously, and the dual-scale repair drives
the two $p_0$ values toward a geometric mean that eliminates the scale
separation~\cite{loutey_paper8}.

This theorem is a guardrail in the project's sense:\ it rules out an
entire class of single-center, shared-exponent molecular encodings
(nested hyperspherical, unified-basis, single-$p_0$ Sturmian CI), and
the resolution it forces---reintroduce the $R$-dependent spectrum by
moving to the prolate spheroidal coordinates where the diatomic
actually separates---is the natural-geometry principle in action.

\subsection{The prolate spheroidal lattice for $\mathrm{H}_2^+$}

Paper~11~\cite{loutey_paper11} implements that resolution.  The
two-center one-electron problem separates in prolate spheroidal
coordinates $(\xi, \eta, \phi)$~\cite{flammer1957}, and the separated equations
discretize on a prolate spheroidal lattice with a spectral angular
solver (Legendre basis) and a spectral radial solver~\cite{trefethen2000}, with Brent
root-finding in the separation parameter $c^2 = -R^2 E_{\mathrm{elec}}/2$
to enforce self-consistency.  The method has zero free parameters:\ all
inputs are nuclear charges, the internuclear distance, and the basis
resolution.  With the spectral Laguerre radial solver at
$n_{\mathrm{basis}} = 20$, the $\mathrm{H}_2^+$ ground-state energy is
$E_{\min} = -0.6026$~Ha against the exact $-0.6026$~Ha~\cite{bates1953}---a $0.0002\%$
error, the most accurate single result in the chemistry
arc---with the correct dissociation to $\mathrm{H}+p$ as
$R \to \infty$~\cite{loutey_paper11}.  The spectral solver achieves a
$250\times$ dimension reduction, a $270\times$ speedup, and a
$5000\times$ accuracy improvement over a finite-difference baseline
whose $1.01\%$ error is shown to be a discretization artifact, not an
intrinsic limit of the basis~\cite{loutey_paper11}.

The algebraic structure here is the first concrete instance of the
arc's central theme.  For $\sigma$ states ($m=0$), the spectral
Laguerre matrix elements are computed entirely from three-term
recurrence relations on Laguerre moment matrices, eliminating numerical
quadrature from the radial solver completely (the overlap is exact,
$S_{mn} = \delta_{mn}/2\alpha$, and the kinetic and potential matrices
are finite sums of moments, reproduced to $\sim\!10^{-12}$ against
quadrature)~\cite{loutey_paper11}.  For $\pi$ and $\delta$ states
($m \neq 0$) an associated Laguerre basis absorbs the centrifugal $1/x$
singularity via a partial-fraction decomposition and a Stieltjes
integral recurrence, reducing the non-algebraic content of the entire
solver to a \emph{single} transcendental seed, the exponential-integral
value $e^a E_1(a)$; all higher moments follow from the rational
recurrence $S_0^{(k)}(a) = \Gamma(k) - a\,S_0^{(k-1)}(a)$ for $k \ge
1$~\cite{loutey_paper11}.  This is the exchange-constant taxonomy of the
foundations arc made concrete:\ the irreducible transcendental content
of the diatomic radial problem is one seed in the \emph{composition}
tier, tagged to a specific integral, rather than diffuse numerical
error~\cite{loutey_paper18}.

%% ========================================================================
\section{Level 3: The Two-Electron Atom and the Cusp}
\label{sec:level3}
%% ========================================================================

\subsection{Algebraic $V_{ee}$ on the prolate spheroidal lattice}

Before leaving the diatomic geometry, Paper~12~\cite{loutey_paper12}
adds the second electron and confronts the electron--electron
repulsion $\langle \phi_i | 1/r_{12} | \phi_j \rangle$ over the
six-dimensional two-particle configuration space.  Standard numerical
quadrature on a grid converges slowly because of the Coulomb
singularity at $r_{12} = 0$ and saturates near $80\%$ of the exact
$\mathrm{H}_2$ dissociation energy regardless of basis size.  The
algebraic replacement is the Neumann expansion~\cite{neumann1878} of $1/r_{12}$ in prolate
spheroidal coordinates, which decomposes each six-dimensional integral
into a finite sum of products of one-dimensional auxiliary integrals
computable by recurrence relations:\ the angular moments $C_l(q)$, the
exponential moments $A_n(\alpha)$, and the first- and second-kind
Legendre moments all satisfy exact three-term recurrences, so the
$V_{ee}$ matrix is exact within the Neumann truncation order with no
quadrature grids and no fitted parameters~\cite{loutey_paper12}.  A
selection rule $C_l(q) = 0$ for $q < l$ or $(q-l)$ odd truncates the
effective sum automatically.

On $\mathrm{H}_2$ at $R = 1.4011$~bohr the algebraic Neumann $V_{ee}$
reaches $92.4\%$ of the exact $D_e$ with $72$ basis functions, against
$80.1\%$ for numerical $V_{ee}$ at the same basis size---a uniform
$12$--$20$ percentage-point improvement at every basis size from $1$ to
$72$ functions, and the Neumann result plateaus, confirming that the
$V_{ee}$ \emph{integration} error has been
eliminated~\cite{loutey_paper12}.  The paper then diagnoses the
remaining $7.6\%$ gap honestly:\ it is a one-electron basis
completeness limit, not an integration error.  The electron--electron
cusp requires the non-analytic terms $r_{12}^{1/2}$ and
$r_{12}\ln r_{12}$ (the Kato/Fock structure~\cite{kato1957}) that no polynomial prolate
spheroidal basis can represent.  This is the framework's embedding-tier
transcendental:\ the cusp is two-dimensional in the angular pair
coordinate and is the canonical example of an exchange constant that no
single algebraic insertion reproduces~\cite{loutey_paper18}.  The
diagnosis identifies the next natural geometry---hyperspherical
coordinates, where the coalescence is a boundary condition rather than
a coordinate singularity.

\subsection{The hyperspherical lattice for helium}

Paper~13~\cite{loutey_paper13} introduces that geometry.  In
hyperspherical coordinates~\cite{aquilanti1986, avery1989} the He problem becomes a coupled-channel
expansion in which the angular eigenvalue problem depends parametrically
on the hyperradius~$R$, so the channel functions vary with~$R$; this
coupled-channel structure has the form of a fiber bundle and is the
mathematical prototype for electron--nuclear coupling in molecules.
The electron--electron cusp at $\alpha = \pi/4$, $\theta_{12}=0$ becomes
a boundary condition rather than a singularity~\cite{loutey_paper13}.

The accuracy results here are reported at carefully separated tiers, and
the synthesis preserves the distinction.  The original single-channel
adiabatic solver gives $E = -2.9052$~Ha, nominally $0.05\%$ from the
exact $-2.9037$~Ha, but it is \emph{non-variational}:\ it lies below the
exact energy through a fortuitous cancellation of adiabatic and
finite-difference error, and the paper labels it as
such~\cite{loutey_paper13}.  The properly variational results come from
a two-dimensional variational solver that treats the hyperradius and
hyperangle simultaneously, avoiding the adiabatic approximation:\
$0.022\%$ raw at $l_{\max}=7$ (a genuine variational upper bound) and
$0.004\%$ with a Schwartz cusp correction at $l_{\max}=4$ (a
non-variational extrapolation that lies marginally below the exact
energy)~\cite{loutey_paper13}.  A separate, fully zero-parameter route
is the graph-native configuration interaction, which uses the graph
Laplacian as the one-body operator with exact rational Slater
integrals and reaches $0.19\%$ at $n_{\max}=7$ with no free
parameters~\cite{loutey_paper13}.  Each route exercises a different
part of the algebraic registry:\ the angular sector is exact throughout
(Gaunt integrals $G(l,k,l') = 2\,(l\,k\,l'\,|\,0\,0\,0)^2$, $SO(6)$
Casimir eigenvalues $\nu(\nu+4)/2$, closed-form perturbation
coefficients), while the adiabatic potential curves $U_{\mathrm{eff}}(R)$
and the parametric angular eigenvalues $\mu(R)$ remain point-by-point
diagonalizations---numerical content named explicitly as a current,
not fundamental, feature~\cite{loutey_paper13, loutey_paper18}.

Paper~13 also documents an instructive negative result:\ in the
\emph{adiabatic} approximation, increasing $l_{\max}$ beyond $0$
worsens the energy, because the finite-difference treatment of the
boundary singularities $l(l+1)/\cos^2\alpha$ and $l(l+1)/\sin^2\alpha$
degrades faster than the partial-wave content improves; the adiabatic
channel truncation has an intrinsic floor near $0.19\%$--$0.20\%$.  The
2D variational solver is what breaks that floor~\cite{loutey_paper13}.
The same fiber-bundle structure, coupled to a nuclear lattice, delivers
an end-to-end demonstration:\ ab initio rovibrational constants for
$\mathrm{H}_2$ with zero spectroscopic input, the fundamental
$\nu_{01} = 4157~\mathrm{cm}^{-1}$ ($0.1\%$) and the rotational
constant $B_e = 59.5~\mathrm{cm}^{-1}$ ($2.2\%$)~\cite{loutey_paper13}.

%% ========================================================================
\section{Level 4: The Two-Center Two-Electron Molecule}
\label{sec:level4}
%% ========================================================================

The two-center two-electron molecule $\mathrm{H}_2$~\cite{james1933} combines all five
Coulomb singularities and had resisted a single natural coordinate
system.  Paper~15~\cite{loutey_paper15} introduces molecule-frame
hyperspherical coordinates $(R_e, \alpha, \theta_1, \theta_2, \Phi)$
parameterized by the dimensionless ratio $\rho = R/(2R_e)$, where $R$
is the internuclear distance and $R_e$ the electronic hyperradius.  The
electron--electron cusp appears at $\alpha = \pi/4$, $\theta_{12} = 0$
\emph{independent of molecular geometry}---a structural advantage over
prolate spheroidal coordinates, where the cusp is a coordinate
singularity---and the angular eigenvalue problem depends on $(R, R_e)$
only through $\rho$, collapsing a two-dimensional parameter sweep to one
dimension~\cite{loutey_paper15}.

A coupled-channel expansion in body-frame partial waves $(l_1, m_1,
l_2, m_2)$ with the gerade constraint $l_1 + l_2$ even and $m_1 + m_2 =
0$, solved with the same two-dimensional $(R_e, \alpha)$ variational
solver that avoids the adiabatic approximation, recovers $96.0\%$ of
the exact $D_e$~\cite{kolos1968} at $l_{\max}=6$ with $\sigma$ and $\pi$ channels ($61$
channels; a complete-basis extrapolation reaches
$\sim\!97\%$)~\cite{loutey_paper15}.  This exceeds the $92.4\%$ that
prolate spheroidal CI with algebraic Neumann integrals achieved
(Paper~12), confirming the cusp-resolution advantage of the
hyperspherical geometry.  The framework extends to heteronuclear
diatomics:\ for $\mathrm{HeH}^+$ a charge-center origin shift and
removal of the gerade constraint recover $93.1\%$ of the exact $D_e$ at
$l_{\max}=4$ with $\sigma+\pi$ channels in the \emph{adiabatic}
treatment~\cite{loutey_paper15}; this is a non-variational over-binding
figure---the consistent two-dimensional variational treatment recovers
only $\sim\!5\%$ of $D_e$, so this is not parity with the
$\mathrm{H}_2$ result.

The algebraic content reaches its cleanest statement here.  The nuclear
potential coupling is a split-region Legendre expansion whose sum over
Gaunt integrals terminates \emph{exactly} at $k = l' + l$ by the $3j$
triangle inequality $|l' - l| \le k \le l' + l$---no truncation
parameter, exact to machine precision in all
regimes~\cite{loutey_paper15}.  This is the arc's prime example of a
continuous integration replaced by a finite algebraic sum once the
right structure is identified.  What remains numerical is named:\ the
$Z_{\mathrm{eff}}$ screening correction (Gauss--Legendre quadrature),
and the angular eigenvalue sweep at each $\rho$, which is a
point-by-point diagonalization because the split-region structure is
piecewise smooth~\cite{loutey_paper15, loutey_paper18}.

Paper~15 is careful about variational hygiene.  The adiabatic
approximation systematically \emph{over}-estimates $D_e$ by roughly
$11\%$ relative to the 2D variational solver, and at $l_{\max}=4$ with
$m_{\max}=1$ the adiabatic $D_e$ exceeds the exact value, violating the
variational principle; the 2D result is the correct bound, and the
$\sim\!10$ percentage-point gap between them is the kinetic-energy cost
of non-adiabatic coupling~\cite{loutey_paper15}.  The convergence also
shows an even--odd staircase:\ even-$l_{\max}$ increments add large
gains (the gerade constraint forces direct nuclear coupling only at even
$l_{\max}$) while odd increments are negligible---a selection-rule
signature, not a numerical accident~\cite{loutey_paper15}.

%% ========================================================================
\section{Level 5: Composed Geometries and the Accuracy Ceiling}
\label{sec:composed}
%% ========================================================================

\subsection{The composition principle}

Systems with core electrons resist a single natural geometry because
core and valence electrons occupy different length scales, and---as
noted in Sec.~\ref{sec:intro}---the hierarchy does not collapse
downward (Level-4 coordinates are undefined for one
electron)~\cite{loutey_paper17}.  Paper~17~\cite{loutey_paper17}
introduces \emph{composed} natural geometries:\ fiber bundles in which
each electron group inhabits its own natural coordinate system, coupled
by screening functions and ab initio Phillips--Kleinman (PK)
pseudopotentials~\cite{phillips_kleinman1959} derived entirely from the core wavefunction.  For LiH,
the $\mathrm{Li}^+\,1s^2$ core is solved by the Level-3 hyperspherical
method, yielding a screening function $Z_{\mathrm{eff}}(r)$ and the PK
parameters; the two valence electrons are solved by the Level-4
molecule-frame method with $Z_{\mathrm{eff}}$ injected; and the
potential-energy surface is scanned by a $\rho$-collapse adiabatic
method.  The construction uses \emph{zero} experimental molecular
data---all PK inputs are atomic (the core wavefunction, the ionization
energy, the nuclear charges)~\cite{loutey_paper17}.

\subsection{Results and the PK accuracy ceiling}

The composed LiH potential-energy surface gives $R_{\mathrm{eq}} =
3.21$~bohr at $l_{\max}=2$ ($6.4\%$ error versus the experimental
$3.015$~bohr~\cite{huber1979}, improved from $17\%$ in the LCAO baseline) and
$\omega_e = 1471~\mathrm{cm}^{-1}$ ($4.6\%$); an $l$-dependent PK
variant, which restricts the projector to the $l=0$ core channels,
reduces the bond-length error to $5.3\%$~\cite{loutey_paper17}.  The
same parameter-free formula transfers to $\mathrm{BeH}_2$ and
$\mathrm{H}_2\mathrm{O}$:\ with full one-reduced-density-matrix
inter-fiber exchange, $\mathrm{BeH}_2$ reaches $11.7\%$
bond-length error (matching a single-parameter fitted exchange model at
zero free parameters), and the uncoupled $\mathrm{H}_2\mathrm{O}$
construction reaches $19.4\%$ (improved from $26\%$ by the v3.x PK-consistency fix) with a correct qualitative
surface~\cite{loutey_paper17}.

Two findings establish that the PK pseudopotential is the
\emph{characterized accuracy ceiling} of the composed architecture, not
a tunable knob.  First, the PK barrier is essential, not cosmetic:\
removing it collapses $R_{\mathrm{eq}}$ to $0.88$~bohr, so PK is what
creates the bonding-versus-Pauli-repulsion balance that produces an
equilibrium at all~\cite{loutey_paper17}.  Second, the residual error
is structural.  Across all four solver$\times$PK combinations the
bond-length error saturates near $5$--$6\%$ at the $l_{\max}=2$
operating point, and the $l_{\max}$ drift (the systematic outward
creep of $R_{\mathrm{eq}}$ with partial-wave order) is not removed by
the 2D variational solver, by the $l$-dependent PK, or by an
$R$-dependent PK whose reference length cannot be derived from atomic
quantities without circular reasoning~\cite{loutey_paper17}.  The paper
attributes this to a deep correlation wall---labeled W1e, the composed
architecture's instance of the framework's multi-focal composition wall
pattern---and documents that single-mechanism repairs each saturate
well below closure:\ a bonding-orbital PK extension at a $43\%$ ceiling,
explicit-core Hartree pressure at $\sim\!26\%$, basis enlargement at
$\sim\!10\%$, with Schmidt orthogonalization and core correlation each
contributing negligibly~\cite{loutey_paper17}.

\subsection{Equilibrium without PK: the Level-4N control}

That PK is a \emph{cost of factorization} rather than a missing physics
term is established by the Level-4N control, the exact $N$-electron
generalization in which all four LiH electrons share one
$S^{11}$ Hilbert space under $SO(12)$ with $S_4$
antisymmetry~\cite{loutey_paper17}.  At $l_{\max}=2$ ($12$ $S_4$
channels, angular dimension $2592$) the Level-4N surface has a genuine
equilibrium near $R_{\mathrm{eq}} \approx 1.0$--$1.1$~bohr ($\approx
64\%$ error) with \emph{no} PK---demonstrating that binding exists
without the pseudopotential, and that the $l_{\max}=2$ $S_4\,[2,2]$
angular basis is genuinely insufficient for LiH binding rather than the
encoding being wrong.  The composed geometry trades this exact
inter-group antisymmetry for a $144\times$ angular compression (cost
scaling $l_{\max}^2$ versus $l_{\max}^{11}$), which is what makes the
composed architecture the production operating point despite the
$5$--$6\%$ ceiling~\cite{loutey_paper17}.  The reason the trade is
forced and not a matter of convenience is a structural obstruction:\
exact antisymmetry requires a shared coordinate system, while
composition factorizes the wavefunction into incompatible coordinate
systems---so the three tested PK alternatives (node exclusion, a
fiber-bundle Berry connection, spectral orthogonality) each fail for the
same reason~\cite{loutey_paper17}.

%% ========================================================================
\section{The PK-Free Alternative and the FCI Encodings}
\label{sec:pkfree_fci}
%% ========================================================================

\subsection{Balanced coupled composition (Paper 19)}

Paper~19~\cite{loutey_paper19} pursues the natural follow-on:\ can the
PK pseudopotential be replaced by explicit cross-block physics?  The
balanced coupled Hamiltonian adds cross-center nuclear attraction
$V_{ne}$ via a multipole expansion with exact Gaunt termination
(the multipole sum terminates at $L_{\max} = 2 l_{\max}$ by the same
$3j$ selection rules, with no truncation error), restoring the
cross-block one-body coupling that the composed architecture
lacks; these missing one-body terms are the Shibuya--Wulfman integrals of Coulomb Sturmian theory~\cite{shibuya1965, loutey_paper19}.  The result is the PK-free alternative for
single-point energies at fixed geometry:\ LiH four-electron FCI reaches
$0.20\%$ energy error at $n_{\max}=3$ (down from $1.8\%$ at $n_{\max}=2$),
using analytical incomplete-gamma cross-center integrals at machine
precision~\cite{loutey_paper19}.  The structural cost is geometric
rather than energetic:\ the equilibrium bond length still drifts to
$R_{\mathrm{eq}} = 3.28$~bohr ($\approx 8.8\%$ error), though the drift
per $n_{\max}$ step is about three times smaller than the PK
$l_{\max}$ drift~\cite{loutey_paper19}.  The honest reading is that
energy converges excellently while geometry drifts, making the balanced
coupled Hamiltonian optimal for single-point quantum simulation at
fixed geometries rather than for full potential-energy
surfaces~\cite{loutey_paper19}.  Paper~19 also records a clean
negative:\ cross-block electron--electron integrals \emph{without} the
matching cross-center $V_{ne}$ are energetically unbalanced ($29\%$
error, worse than PK), and a harmonic-driven orthogonalization wall
blocks second-row hydrides such as NaH, which PK cross-center coupling
reduces but does not close~\cite{loutey_paper19}.

\subsection{Graph-native FCI for atoms}

The atomic FCI paper~\cite{loutey_fci_atoms} solves multi-electron
atoms directly on the sparse $S^3$ graph Hamiltonian using
Slater--Condon rules~\cite{Slater1960, CondonShortley1935} with exact $R^k$ radial integrals and Gaunt
coefficients, with Slater determinants indexed as sorted tuples of
spin-orbital labels (relational lattice indexing).  Helium reaches
$0.35\%$ at $n_{\max}=5$ (grid-based) and $0.19\%$ at $n_{\max}=7$ with
exact rational Slater integrals; lithium $1.07\%$ (at $n_{\max}=5$)
and beryllium $0.90\%$ (at $n_{\max}=4$, with $487{,}635$
determinants)~\cite{loutey_fci_atoms}.  The Hamiltonian is sparse
(fill fraction $0.68\%$ for He at $n_{\max}=5$, $0.12\%$ for Li at
$n_{\max}=4$) with a sub-quadratic nonzero-count growth exponent
$\alpha \approx 1.4$--$1.5$.  The paper documents an instructive
variational subtlety:\ a hybrid one-body operator that mixes graph
Laplacian off-diagonals (carrying $\kappa = -1/16$) improves He but
violates the variational principle for three-electron lithium, so the
exact hydrogenic one-body operator is required for monotonic
convergence~\cite{loutey_fci_atoms}.  The dominant remaining error for
heteronuclear LiH at $n_{\max}=3$ is basis truncation, not integral
incompleteness---the basis-set superposition error there exceeds the
experimental binding energy, so the raw binding number is reported with
its counterpoise correction~\cite{boys_bernardi} and its honest caveat~\cite{loutey_fci_atoms}.

\subsection{The graph-concatenation theorem (a second guardrail)}

The molecular FCI paper~\cite{loutey_fci_mol} supplies the arc's second
proven negative theorem.  A topological LCAO architecture that builds a
molecular graph by concatenating atom-centered hydrogenic lattices
(via Mulliken cross-nuclear attraction, Slater--Condon repulsion, and
overlap bridges) is shown to have an intra-atomic graph-Laplacian
kinetic energy that is \emph{completely $R$-independent} by
construction:\ the atomic lattice structure does not change as the
nuclei approach.  The balanced Hamiltonian therefore produces a
monotonically attractive potential-energy surface with no equilibrium
geometry---the binding energy decreases smoothly from $0.464$~Ha at
$R=1.5$~bohr to $0.042$~Ha at $R=5.0$~bohr, with a virial ratio far
from unity~\cite{loutey_fci_mol}.  The paper documents a $29$-version
diagnostic arc that falsifies basis-set superposition error as the
geometric driver, rules out orbital-exponent relaxation, and---citing
the Sturmian structural theorem of Paper~8---disproves every
single-center Sturmian and bond-sphere alternative as a repair; a
single-parameter Fock-weighted kinetic correction partially restores an
equilibrium ($R_{\mathrm{eq}} \approx 3.8$~bohr) but cannot fit
$R_{\mathrm{eq}}$ and $D_e$ simultaneously~\cite{loutey_fci_mol,
loutey_paper8}.

Together the two negative theorems---Sturmian structural (shared
$SO(4)$ Wigner-$D$ congruence $\Rightarrow$ $R$-independent eigenvalues,
Paper~8) and graph
concatenation ($R$-independent kinetic energy $\Rightarrow$ no
equilibrium, the molecular FCI paper)---are guardrails:\ they explain
\emph{why} the natural-geometry hierarchy is necessary.  A molecule
must be discretized in coordinates where the separation introduces
$R$-dependence (prolate spheroidal, molecule-frame hyperspherical, or
composed), because the two obvious shortcuts (one shared-exponent
single-center basis; one concatenated graph) are provably incapable of
binding.

%% ========================================================================
\section{The Algebraic-First Thread}
\label{sec:algebraic}
%% ========================================================================

A single discipline runs through every level of the hierarchy:\ replace
continuous integration with algebraic evaluation wherever the structure
allows, and name what remains.  This is the chemistry-arc instance of
the framework's prime directive (the discrete angular/selection-rule
structure is exact at all levels and must never be approximated; the
radial/parametric content is level-dependent and a legitimate target
for numerical improvement).  Table~\ref{tab:algebraic} catalogues the
state of the algebraic registry across the arc.

\begin{table*}[t]
\centering
\caption{The algebraic-first registry across the chemistry arc.
\emph{Algebraic} = closed-form or exact-recurrence from quantum-number
labels with no quadrature; \emph{algebraic (one seed)} = exact up to a
single tagged transcendental; \emph{numerical} = currently requires
quadrature or point-by-point diagonalization, named as a present (not
fundamental) feature.}
\label{tab:algebraic}
\begin{tabular}{l l l}
\toprule
Level / object & Status & Mechanism (paper) \\
\midrule
Angular coupling (all levels) & algebraic & Gaunt $/$ Wigner $3j$; selection rules \cite{loutey_paper13, loutey_paper15, biedenharn1981} \\
$\mathrm{H}_2^+$ radial, $\sigma$ ($m=0$) & algebraic & three-term Laguerre recurrence \cite{loutey_paper11} \\
$\mathrm{H}_2^+$ radial, $\pi/\delta$ ($m\neq0$) & algebraic (one seed) & assoc.\ Laguerre $+$ seed $e^a E_1(a)$ \cite{loutey_paper11} \\
Prolate spheroidal $V_{ee}$ & algebraic & Neumann expansion, recurrence moments \cite{loutey_paper12} \\
He hyperradial $S$, $K$ matrices & algebraic & three-term Laguerre recurrence \cite{loutey_paper13} \\
Nuclear coupling (Level 4) & algebraic & split-region Legendre, $3j$-triangle termination \cite{loutey_paper15} \\
Cross-center $V_{ne}$ (balanced) & algebraic & multipole, exact Gaunt termination $L_{\max}=2l_{\max}$ \cite{loutey_paper19} \\
Slater $F^k$ integrals (graph-native) & algebraic & exact rational $S^3$ node overlap \cite{loutey_paper13, loutey_fci_atoms} \\
$Z_{\mathrm{eff}}$ screening & numerical & Gauss--Legendre quadrature \cite{loutey_paper15, loutey_paper17} \\
Adiabatic potential $U(R)$, $\mu(R)$ & numerical & point-by-point diagonalization \cite{loutey_paper13} \\
Hyperradial coupling (Levels 3--4) & numerical & FD $/$ spectral radial solve \cite{loutey_paper13, loutey_paper15} \\
\bottomrule
\end{tabular}
\end{table*}

The pattern is consistent with the foundations-arc taxonomy of
Paper~18~\cite{loutey_paper18}.  The angular sector is intrinsic
(rational, $\pi$-free).  The diatomic radial seed $e^a E_1(a)$ and the
Slater-integral seeds are composition-tier transcendentals tied to a
specific integral.  The electron--electron cusp is the embedding-tier
exchange constant.  And the items still in the numerical column---the
$Z_{\mathrm{eff}}$ screening, the adiabatic curves, the hyperradial
coupling---are precisely the places where an algebraic structure has
not yet been found, which the project records as an ongoing goal rather
than a deficiency:\ the Neumann expansion (Paper~12) and the
split-region Legendre termination (Paper~15) are existence proofs that
quadrature walls in this arc \emph{can} dissolve once the right
structure is identified~\cite{loutey_paper12, loutey_paper15}.

%% ========================================================================
\section{Positioning: A Research Instrument}
\label{sec:positioning}
%% ========================================================================

The classical-solver investigation summarized in this synthesis is
complete:\ the natural-geometry hierarchy has been built level by level,
its algebraic content catalogued, and its accuracy ceilings
characterized.  The framework's value proposition is best stated against
those ceilings honestly.

\subsection{What the framework offers}

The construction is \emph{zero-parameter}:\ every result in
Table~\ref{tab:hierarchy} takes as input only nuclear charges,
geometry, and a basis resolution---no fitted potentials, no empirical
screening constants, no experimental molecular data (the PK
pseudopotential of Paper~17 is derived entirely from the atomic core
wavefunction)~\cite{loutey_paper11, loutey_paper17}.  Angular-momentum
selection rules are baked into the basis, so the matrix elements that
vanish by symmetry are never computed, and the surviving ones are exact
algebraic quantities.  The Hamiltonians are structurally sparse, with
sub-quadratic nonzero growth~\cite{loutey_fci_atoms}.  These structural
properties---not interpretive claims about the nature of quantum
mechanics---are the framework's actual advantages.

\subsection{Where the headline value lives}

The arc's strongest computational payoff is not the classical
ground-state energies but the qubit Hamiltonians they imply.  The
angular sparsity that this synthesis documents at the integral level is
the structural reason the composed-geometry qubit encodings achieve
favorable Pauli-term scaling for variational quantum simulation; that
result, its scaling exponents, and its molecular benchmarks belong to
the quantum-computing arc (Group~4) and are referenced here, not
re-derived.  This synthesis establishes the chemistry-side foundation of
that payoff:\ the natural-geometry hierarchy, the algebraic-first
integral structure, and the two negative theorems that justify the
composed encoding.

\subsection{The benchmarking discipline}

Where the framework is compared to other methods, the comparison is to
the strongest available baseline, and unfavorable comparisons are stated
plainly.  The composed potential-energy surfaces carry $5$--$19.4\%$
bond-length error, well above production quantum-chemistry accuracy; the
honest reading is that the framework is a computationally principled
\emph{alternative} offering sparsity, scaling, and structural insight,
not a replacement for correlated electronic-structure methods.  The
papers report the proven negatives---the Sturmian and concatenation
theorems, the PK ceiling, the W1e wall, the failed PK alternatives and
cusp treatments---as a first-class part of the corpus, because the
characterization of what natural geometry \emph{cannot} do is as
load-bearing as what it can~\cite{loutey_paper8, loutey_fci_mol,
loutey_paper17}.

%% ========================================================================
\section{Open Questions and Outlook}
\label{sec:open}
%% ========================================================================

Three structural questions remain open at the close of the
classical-solver arc.

\paragraph{The composed accuracy ceiling (W1e).}  The composed
architecture saturates near $5$--$6\%$ bond-length error at its
$l_{\max}=2$ operating point, with the residual diagnosed as a
multi-determinant correlation wall whose single-mechanism repairs each
close only a fraction~\cite{loutey_paper17}.  Whether a composition
rule exists that closes W1e without sacrificing the $144\times$ angular
compression---i.e.\ without reverting to the exact-but-expensive
Level-4N encoding---is the central open chemistry question.  The
diagnostic record (the failed PK alternatives, the structural
obstruction that exact antisymmetry requires a shared coordinate
system) suggests the answer is tied to the framework's
sparsity-versus-binding trade-off rather than to a missing integral.

\paragraph{Algebraic replacement of the remaining quadrature.}  The
numerical column of Table~\ref{tab:algebraic}---$Z_{\mathrm{eff}}$
screening, the adiabatic potential curves, the hyperradial
coupling---is the standing target of the algebraic-first program.  The
existence proofs (Neumann $V_{ee}$, split-region Legendre termination)
show such replacements are achievable in principle; whether the
adiabatic curves $\mu(R)$, which Paper~13 records as genuinely
transcendental in $R$, admit an algebraic-implicit form (defined by a
polynomial equation over a known coefficient ring) is open.

\paragraph{Second-row and polyatomic extension.}  The balanced coupled
Hamiltonian binds first-row hydrides but hits an orthogonalization wall
for second-row systems such as NaH~\cite{loutey_paper19}, and the
polyatomic $\mathrm{H}_2\mathrm{O}$ construction validates only
bond--bond coupling, with bond--lone and lone--lone couplings
unphysical at the present level~\cite{loutey_paper17}.  Both are named
boundaries of the current encoding rather than refutations of the
hierarchy.

None of these open questions destabilizes the structural results of the
arc.  The natural-geometry hierarchy is established level by level; the
algebraic-first content is catalogued; the two negative theorems are
proven; and the accuracy ceilings are characterized.  The open
questions are frontier questions about how far the framework's
sparsity can be pushed before its accuracy ceiling binds, not
foundational concerns.

%% ========================================================================
\section{Conclusion}
\label{sec:conclusion}
%% ========================================================================

The Group~2 quantum-chemistry arc of GeoVac is the systematic
construction of the natural-geometry hierarchy for electronic
structure.  Each level is the coordinate system in which a given system
separates:\ $S^3$ for atoms~\cite{loutey_paper7}, the prolate spheroid
for $\mathrm{H}_2^+$ ($0.0002\%$)~\cite{loutey_paper11}, hyperspherical
coordinates for He ($0.022\%$ variational raw / $0.004\%$ non-variational
cusp, $0.19\%$ graph-native)~\cite{loutey_paper13}, molecule-frame hyperspherical
coordinates for $\mathrm{H}_2$ ($96.0\%$ of
$D_e$)~\cite{loutey_paper15}, and composed fiber bundles for
core--valence and polyatomic systems ($5.3\%$, $11.7\%$, $19.4\%$
bond-length error for LiH, $\mathrm{BeH}_2$,
$\mathrm{H}_2\mathrm{O}$)~\cite{loutey_paper17}.  The arc is held
together by an algebraic-first discipline---exact angular coupling,
algebraic Neumann $V_{ee}$~\cite{loutey_paper12}, Laguerre recurrences,
and $3j$-terminating Legendre expansions~\cite{loutey_paper15}---and is
bounded by two proven negative theorems:\ a single-center shared-$p_0$
Sturmian basis cannot bind a molecule~\cite{loutey_paper8}, and a
concatenated atom-centered graph has no equilibrium~\cite{loutey_fci_mol}.
Those theorems are why the hierarchy is necessary, not optional.

The framework is a zero-parameter, structurally sparse research
instrument whose accuracy ceilings are now characterized and whose
strongest downstream payoff---qubit-Hamiltonian sparsity for quantum
simulation---lives in the quantum-computing arc.  The discrete graph
and the continuous Laplace--Beltrami operator are dual descriptions
connected by proven conformal maps; the framework's concrete
advantages do not depend on which description one regards as
fundamental, and that interpretive question is left open.

\section*{Acknowledgments}

This synthesis paper consolidates work across the Group~2
quantum-chemistry arc of the GeoVac project.  The author thanks the
broader open-source computational quantum-chemistry community for the
libraries (NumPy, SciPy, sympy, mpmath, PySCF) that make this work
possible.  The framework's papers and code are publicly archived on
Zenodo with permanent DOIs; this synthesis is a project-archive
document and is not formatted for journal submission.

%================================================================
\begin{thebibliography}{99}
%================================================================

%---- External / classical references ----

\bibitem{fock1935}
V.~Fock,
``Zur Theorie des Wasserstoffatoms,''
Z.\ Phys.\ \textbf{98}, 145 (1935).

\bibitem{bates1953}
D.~R.~Bates, K.~Ledsham, and A.~L.~Stewart,
``Wave functions of the hydrogen molecular ion,''
Philos.\ Trans.\ R.\ Soc.\ Lond.\ A \textbf{246}, 215 (1953).

\bibitem{flammer1957}
C.~Flammer,
\textit{Spheroidal Wave Functions}
(Stanford University Press, Stanford, 1957).

\bibitem{neumann1878}
C.~Neumann,
\textit{\"Uber die Entwicklung einer Function nach den Kugelfunctionen}
(Teubner, Leipzig, 1878).

\bibitem{james1933}
H.~M.~James and A.~S.~Coolidge,
``The ground state of the hydrogen molecule,''
J.\ Chem.\ Phys.\ \textbf{1}, 825 (1933).

\bibitem{kolos1968}
W.~Ko\l{}os and L.~Wolniewicz,
``Improved theoretical ground-state energy of the hydrogen molecule,''
J.\ Chem.\ Phys.\ \textbf{49}, 404 (1968).

\bibitem{kato1957}
T.~Kato,
``On the eigenfunctions of many-particle systems in quantum mechanics,''
Commun.\ Pure Appl.\ Math.\ \textbf{10}, 151 (1957).

\bibitem{aquilanti1986}
V.~Aquilanti, S.~Cavalli, and G.~Grossi,
``Hyperspherical coordinates for molecular dynamics by the method of trees,''
J.\ Chem.\ Phys.\ \textbf{85}, 1362 (1986).

\bibitem{avery1989}
J.~Avery,
\textit{Hyperspherical Harmonics: Applications in Quantum Theory}
(Kluwer Academic, Dordrecht, 1989).

\bibitem{avery2004}
J.~Avery and J.~Avery,
``Generalized Sturmians and atomic spectra,''
Adv.\ Quantum Chem.\ \textbf{47}, 159 (2004).

\bibitem{shibuya1965}
T.~Shibuya and C.~E.~Wulfman,
``Molecular orbitals in momentum space,''
Proc.\ R.\ Soc.\ Lond.\ A \textbf{286}, 376 (1965).

\bibitem{biedenharn1981}
L.~C.~Biedenharn and J.~D.~Louck,
\textit{Angular Momentum in Quantum Physics}
(Addison-Wesley, Reading, MA, 1981).

\bibitem{CondonShortley1935}
E.~U.~Condon and G.~H.~Shortley,
\textit{The Theory of Atomic Spectra}
(Cambridge University Press, Cambridge, 1935).

\bibitem{Slater1960}
J.~C.~Slater,
\textit{Quantum Theory of Atomic Structure}, Vols.~I and~II
(McGraw-Hill, New York, 1960).

\bibitem{boys_bernardi}
S.~F.~Boys and F.~Bernardi,
``The calculation of small molecular interactions by the differences of separate total energies,''
Mol.\ Phys.\ \textbf{19}, 553 (1970).

\bibitem{huber1979}
K.~P.~Huber and G.~Herzberg,
\textit{Molecular Spectra and Molecular Structure IV: Constants of Diatomic Molecules}
(Van Nostrand Reinhold, New York, 1979).

\bibitem{phillips_kleinman1959}
J.~C.~Phillips and L.~Kleinman,
``New method for calculating wave functions in crystals and molecules,''
Phys.\ Rev.\ \textbf{116}, 287 (1959).

\bibitem{trefethen2000}
L.~N.~Trefethen,
\textit{Spectral Methods in MATLAB}
(SIAM, Philadelphia, 2000).

%---- GeoVac papers ----

\bibitem{loutey_paper0}
J.~Loutey,
``Angular Momentum as Information Geometry: Isotropic Packing and the Universal Angular Cross-Section,''
GeoVac Paper~0 (2026).

\bibitem{loutey_paper1}
J.~Loutey,
``The Geometric Atom: Atomic Structure as a Packing Problem,''
GeoVac Paper~1 (2026).

\bibitem{loutey_paper7}
J.~Loutey,
``The Dimensionless Vacuum: Recovering the Schr\"odinger Equation from Scale-Invariant Graph Topology,''
GeoVac Paper~7 (2026).

\bibitem{loutey_paper8}
J.~Loutey,
``Bond Sphere: Diatomic Molecules as $S^3$ and their Natural Sturmian Basis,''
GeoVac Papers~8--9 (2026).

\bibitem{loutey_paper11}
J.~Loutey,
``The Molecular Fock Projection: Diatomic Molecules as Prolate Spheroidal Lattices,''
GeoVac Paper~11 (2026).

\bibitem{loutey_paper12}
J.~Loutey,
``Algebraic Two-Electron Integrals on the Prolate Spheroidal Lattice,''
GeoVac Paper~12 (2026).

\bibitem{loutey_paper13}
J.~Loutey,
``The Hyperspherical Lattice: Two-Electron Atoms as Coupled Channel Graphs,''
GeoVac Paper~13 (2026).

\bibitem{loutey_paper15}
J.~Loutey,
``The Level~4 Natural Geometry: Two-Center Two-Electron Molecules in Molecule-Frame Hyperspherical Coordinates,''
GeoVac Paper~15 (2026).

\bibitem{loutey_paper17}
J.~Loutey,
``Composed Natural Geometries for Core--Valence Diatomics: LiH and BeH$^+$ from Discrete Graph Topology,''
GeoVac Paper~17 (2026).

\bibitem{loutey_paper18}
J.~Loutey,
``Spectral-Geometric Exchange Constants: Projecting Discrete Spectra onto Continuous Manifolds,''
GeoVac Paper~18 (2026).

\bibitem{loutey_paper19}
J.~Loutey,
``Toward PK-Free Molecular Hamiltonians: Two-Center Sturmian Integrals and Coupled Composition,''
GeoVac Paper~19 (2026).

\bibitem{loutey_fci_atoms}
J.~Loutey,
``Full Configuration Interaction on a Sparse Graph Hamiltonian: Multi-electron Atoms via Relational Lattice Indexing,''
GeoVac FCI (atoms) Paper (2026).

\bibitem{loutey_fci_mol}
J.~Loutey,
``Topological Full Configuration Interaction for Heteronuclear Diatomics: LiH as a Benchmark,''
GeoVac FCI (molecules) Paper (2026).

\end{thebibliography}

\end{document}
